% ============================================================================
% Fichier  : partie1/P01_C03_grandes_affaires.tex
% Titre    : Les Grandes Affaires --- Leçons de l'Histoire
% Auteur   : MINKA MI NGUIDJOÏ Thierry Emmanuel
% Version  : 1.0 -- Juin 2026
% Statut   : RÉDIGÉ
%
% Positionnement par rapport au Chapitre 2 :
%   Le Chap. 2 a utilisé les affaires comme preuves de changements de régime.
%   Ce chapitre les dissèque comme laboratoires techniques : anatomie de
%   l'attaque, outils mobilisés, erreurs commises, innovation produite,
%   leçon opérationnelle extractible. Aucune duplication -- complémentarité.
%
% Sources :
%   - 2_grandes_affaires.tex (BTK, Stuxnet, WannaCry -- squelette)
%   - 15_affaire_cyberfinance.tex (CyberFinance Cameroun 2025)
%   - M7_Cas_Pratiques_Integrateurs_PNC.docx (ONAMBELE, TANGUE, MVOM, WOURI)
%   - chap2.pdf (analyses épistémologiques complémentaires)
% ============================================================================

\chapter{Les Grandes Affaires --- Leçons de l'Histoire}
\label{chap:gr-affaires}

\epigraph{%
    « Chaque grande affaire forensique est un laboratoire où s'expérimente
    l'avenir de notre discipline. »%
}{--- D\textsuperscript{r} Henry C. Lee}

\niveauChapitre{Fondamental}{%
    Ce chapitre prolonge le Chapitre~\ref{chap:histoire} en approfondissant
    l'anatomie technique des affaires fondatrices. Aucun prérequis -- les
    concepts techniques sont introduits et définis au fil du texte.%
}

\begin{boxobjectifs}
\begin{itemize}
    \item Analyser l'anatomie technique et judiciaire de sept affaires
          fondatrices pour en extraire les principes forensiques durables.
    \item Identifier, dans chaque affaire, l'\termecle{artefact décisif} ---
          la preuve unique qui a fait basculer le dossier --- et la méthode
          employée pour le découvrir.
    \item Reconnaître les erreurs procédurales récurrentes qui affaiblissent
          les dossiers, afin de ne pas les reproduire.
    \item Appliquer la grille d'analyse CRO à des cas réels pour comprendre
          comment les tensions entre confidentialité, fiabilité et opposabilité
          se manifestent concrètement.
    \item S'approprier les quatre cas camerounais (Opérations ONAMBELE, TANGUE,
          MVOM, WOURI) comme référentiels de pratique locale.
\end{itemize}
\end{boxobjectifs}

\bigskip

% ============================================================================
\section*{Comment lire ce chapitre}
% ============================================================================

Chaque affaire est disséquée selon la même structure~: \textbf{contexte} (ce
qui s'est passé, en bref), \textbf{artefact décisif} (la preuve qui a tout
changé), \textbf{anatomie technique} (comment cette preuve a été trouvée et
pourquoi elle a tenu juridiquement), \textbf{erreurs commises} (ce qui a failli
faire échouer l'investigation), et \textbf{leçon transférable} (ce que vous
devez emporter de cette affaire dans votre pratique quotidienne).

\medskip

Les quatre premières affaires sont internationales et fondatrices. Les quatre
suivantes sont camerounaises --- fictives mais construites sur les patterns
exacts des affaires traitées par la \sigle{DGSN} et l'\sigle{ANTIC} entre
2020 et 2026.

\separateur

% ============================================================================
\section{L'Affaire BTK Killer --- Dennis Rader (2005)}
\label{sec:btk}
% ============================================================================

\subsection{Contexte et artefact décisif}

Dennis Rader, tueur en série surnommé \termecle{BTK} (\textit{Bind, Torture,
Kill}), a assassiné 10~personnes dans la région de Wichita (Kansas) entre
1974 et 1991. Après une longue absence, il reprend contact avec la police en
2004 par lettres, puis par disquette. En janvier~2005, il envoie une disquette
contenant un fichier \texttt{Test.A.rtf}. Les métadonnées extraites par les
enquêteurs du FBI en quelques minutes conduisent à son arrestation en
quatre jours.

\medskip

\textbf{L'artefact décisif}~: les propriétés du fichier \texttt{.rtf} révèlent
le champ \textit{Auteur}~: \texttt{Dennis}, et le champ
\textit{Dernière organisation}~: \texttt{Christ Lutheran Church}. Le
croisement avec les membres actifs de cette paroisse identifie Rader
immédiatement.

\subsection{Anatomie technique}

\begin{boxartefact}{Métadonnées Office --- champs auteur et organisation}
\textbf{Localisation dans Windows~:}
Clic droit sur le fichier $\to$ \texttt{Propriétés} $\to$
onglet \texttt{Détails}. Ou dans Office~: \texttt{Fichier} $\to$
\texttt{Informations} $\to$ \texttt{Propriétés}.

\textbf{Localisation dans Autopsy~:}
\texttt{Results} $\to$ \texttt{Extracted Content} $\to$
\texttt{Document Metadata}.

\textbf{Extraction en ligne de commande~:}
\begin{lstlisting}[language=bash_forensique]
# ExifTool -- extraction complete des metadonnees
exiftool fichier.docx

# Sortie typique
Author              : Dennis
Last Modified By    : Dennis
Organization        : Christ Lutheran Church
Create Date         : 2005:01:15 19:42:33
\end{lstlisting}

\textbf{Principe clé~:} tout document Office créé depuis 1997 porte ces
métadonnées. La plupart des utilisateurs ignorent leur existence.
Elles survivent au renommage, à la copie, au transfert par email.
\end{boxartefact}

\subsection{Erreurs commises --- et ce qu'elles enseignent}

Rader a commis une erreur élémentaire~: utiliser un ordinateur institutionnel
(celui de son église) pour préparer un document destiné à la police, sans
jamais nettoyer les métadonnées. Cette erreur révèle un \textbf{angle mort
cognitif} universel~: les utilisateurs visualisent le contenu de leurs
documents mais ignorent la couche de métadonnées sous-jacente. Pour
l'investigateur, cet angle mort est une opportunité systématique.

\begin{boxalert}
\textbf{Le principe des métadonnées survivantes}~: tout fichier numérique
produit une \emph{deuxième couche d'information} --- les métadonnées --- que
son auteur ne voit pas et oublie donc de supprimer. Cette couche contient
souvent~: l'identité de l'auteur, l'organisation, l'équipement utilisé,
le chemin de fichier sur l'ordinateur source, parfois des coordonnées GPS
(photos), la géolocalisation de l'imprimante, les révisions successives du
document. L'examen systématique des métadonnées est une étape non
optionnelle de toute investigation numérique.
\end{boxalert}

\subsection{Leçon transférable}

\begin{boxresume}
\textbf{BTK --- Ce que vous emportez~:}
\begin{itemize}
    \item Extraire systématiquement les métadonnées de \textbf{tous} les
          fichiers collectés~; ne jamais se limiter au contenu visible.
    \item La commande \lstinline[language=bash_forensique]{exiftool -r /dossier/}
          traite récursivement un répertoire entier.
    \item Un champ auteur vide n'est pas neutre~: il signifie que le fichier
          a été créé avec un outil qui ne l'inscrit pas, ou que les
          métadonnées ont été \emph{délibérément} supprimées (\textit{timestomping}
          ou nettoyage EXIF) --- ce qui est en soi un indice forensique.
\end{itemize}
\end{boxresume}

% ============================================================================
\section{L'Affaire Stuxnet (2010)}
\label{sec:stuxnet}
% ============================================================================

\subsection{Contexte et artefact décisif}

\termecle{Stuxnet} est le premier programme malveillant conçu pour causer des
dommages physiques à une infrastructure industrielle. Découvert en juin~2010
par VirusBlokAda, il cible spécifiquement les centrifugeuses d'enrichissement
d'uranium iraniennes de Natanz. Quatre zero-days simultanés, un code polymorphe
sophistiqué, une précision chirurgicale sur les automates Siemens
\texttt{SIMATIC S7-315/S7-417}~: l'attribution à des services étatiques
(États-Unis et Israël, selon les sources) est largement acceptée.

\medskip

\textbf{L'artefact décisif}~: la présence de \textbf{quatre zero-days
simultanés} dans un seul programme. Un groupe cybercriminel ordinaire
en exploite rarement plus d'un~; en utiliser quatre simultanément signale
des ressources d'État. Ce \textit{marqueur de sophistication} a été
l'argument principal de l'attribution.

\subsection{Anatomie technique}

L'investigation de Stuxnet a fondé la sous-discipline de la
\termecle{cyberforensique industrielle}\index{cyberforensique industrielle}.
Quatre techniques innovantes ont été nécessaires~:

\begin{enumerate}
    \item \textbf{Reverse engineering de code polymorphe~:} Stuxnet modifiait
          son propre code à chaque exécution pour contourner les signatures
          antivirales. L'analyse a requis l'extraction du code depuis la
          mémoire du processus en cours d'exécution --- avant que la
          polymorphie ne l'altère.

    \item \textbf{Sandboxing comportemental~:} l'exécution dans un
          environnement virtualisé isolé (\textit{sandbox}) a permis
          d'observer le comportement du malware sans risquer de propager
          l'infection. Stuxnet vérifiait lui-même s'il tournait dans une
          sandbox~; les analystes ont dû contourner cette détection.

    \item \textbf{Identification des quatre zero-days~:}
          \begin{itemize}
              \item CVE-2010-2568~: vulnérabilité dans le traitement des
                    fichiers \texttt{.lnk} (raccourcis Windows).
              \item CVE-2010-2772~: vulnérabilité dans le \textit{Windows
                    Task Scheduler}.
              \item CVE-2010-2729~: vulnérabilité dans le service
                    \textit{Print Spooler} pour la propagation réseau.
              \item CVE-2010-2743~: vulnérabilité dans le pilote
                    noyau Windows pour l'élévation de privilèges.
          \end{itemize}

    \item \textbf{Forensique des automates industriels (ICS/SCADA)~:}
          les \sigle{PLC} (\textit{Programmable Logic Controllers}) Siemens
          ciblés ne sont pas des PC~; leur forensique exige des outils et
          des compétences spécifiques absents des boîtes à outils classiques.
\end{enumerate}

\begin{boxCRO}
\textbf{Analyse CRO --- Affaire Stuxnet}

$\mathcal{C}$~: l'investigation a nécessité l'accès à des données sensibles
d'États souverains et d'installations nucléaires. Le niveau de
classification a fortement contraint la diffusion des résultats.
\cro{$\downarrow$ très contrainte}{$\uparrow$ robuste}{$\sim$ limitée aux instances spécialisées}

La tension CRO de Stuxnet est atypique~: l'opposabilité n'est pas un
objectif pénal (personne ne sera jugé) mais stratégique (attribution
pour la diplomatie et la dissuasion). Cela illustre que l'investigation
numérique sert des finalités multiples, pas seulement judiciaires.
\end{boxCRO}

\subsection{Leçon transférable}

\begin{boxresume}
\textbf{Stuxnet --- Ce que vous emportez~:}
\begin{itemize}
    \item La sophistication d'un malware est un \textbf{marqueur
          d'attribution}~: zéro à deux zero-days = acteur criminel ou
          hacktiviste~; trois zero-days et plus = ressources d'État.
    \item Les systèmes industriels (ICS/SCADA) sont une frontière en
          expansion de l'investigation numérique~; leur forensique
          nécessite des compétences spécifiques au-delà de l'informatique
          classique.
    \item L'analyse comportementale en sandbox est indispensable pour
          tout malware inconnu~: ne jamais exécuter un binaire suspect
          sur un système de production ou d'analyse non isolé.
\end{itemize}
\end{boxresume}

% ============================================================================
\section{L'Affaire WannaCry (2017)}
\label{sec:wannacry}
% ============================================================================

\subsection{Contexte et artefact décisif}

Le 12~mai~2017, le ransomware \termecle{WannaCry} infecte 300~000~ordinateurs
dans 150~pays en moins de 72~heures, en exploitant la vulnérabilité
\textit{EternalBlue} (CVE-2017-0144) volée à la NSA et publiée par le
groupe Shadow Brokers. Les systèmes Windows non patchés, notamment ceux du
National Health Service britannique (NHS), sont frappés de plein fouet~:
des hôpitaux doivent reporter des opérations chirurgicales non urgentes.

\medskip

\textbf{L'artefact décisif}~: le \textbf{kill switch}. Marcus Hutchins,
chercheur de 22~ans, découvre dans le code de WannaCry une requête vers un
domaine non enregistré (\texttt{iuqerfsodp9ifjaposdfjhgosurijfaewrwergwea.com}).
En enregistrant ce domaine pour 10,69~USD, il active le kill switch et stoppe
la propagation mondiale.

\subsection{Anatomie technique}

Trois techniques d'investigation ont été déterminantes~:

\subsubsection{Analyse du Domain Generation Algorithm (DGA)}

WannaCry ne s'appuie pas sur un DGA classique --- le domaine kill switch
est codé en dur. Mais l'analyse de ce domaine a révélé l'architecture de
contrôle du malware~: la présence d'un kill switch hard-codé est une
signature des opérations de malware à grande échelle (permet à l'auteur de
désactiver le malware en cas d'attribution imminente).

\begin{lstlisting}[language=bash_forensique,
    caption={Extraction du kill switch depuis un binaire WannaCry}]
# Strings -- recherche de domaines et URLs dans le binaire
strings wannacry.exe | grep -E "\.(com|net|org|onion)"

# Sortie attendue :
# iuqerfsodp9ifjaposdfjhgosurijfaewrwergwea.com  <- kill switch
# 2.0.0.10 <- C2 IP candidate
# @WanaDecryptor@  <- marqueur de famille

# Vérification dynamique en sandbox (REMnux)
remnux analyze --timeout 60 wannacry.exe
\end{lstlisting}

\subsubsection{Attribution probabiliste~: la piste Lazarus}

L'attribution à la Corée du Nord (groupe Lazarus) repose sur trois
artefacts corrélés~: similitudes de code avec des malwares Lazarus
précédents (2014~: Sony Pictures), utilisation de la même adresse
Bitcoin pour les paiements de rançons, et présence de commentaires
en coréen dans certaines chaînes de débogage.

\begin{boiteRemarque}{Attribution numérique~: probabilité, pas certitude}
L'attribution d'une attaque numérique à un acteur étatique ou
criminel est un \textbf{jugement probabiliste}, jamais une certitude.
Les bons acteurs imitent les signatures d'autres groupes (\textit{false
flag}). L'investigateur doit présenter ses conclusions d'attribution avec
un niveau de confiance explicite (\textit{ex.}~: « attribution avec
confiance modérée, sur la base de trois indicateurs corrélés »), non
comme des faits établis.
\end{boiteRemarque}

\subsubsection{Forensique de l'impact~: reconstitution de la chronologie}

La réponse à WannaCry a révélé l'importance de la
\termecle{forensique des journaux distribués}\index{forensique!journaux distribués}~:
les logs de 150~pays devaient être corrélés temporellement. La
leçon fondamentale est que les horloges système ne sont pas synchronisées
universellement --- une écart de quelques minutes entre deux systèmes
peut inverser la chronologie d'une infection et fausser l'attribution.

\subsection{Leçon transférable}

\begin{boxresume}
\textbf{WannaCry --- Ce que vous emportez~:}
\begin{itemize}
    \item Analyser systématiquement les requêtes réseau d'un malware~:
          elles révèlent l'architecture de contrôle et parfois un
          mécanisme de désactivation.
    \item L'attribution est une \textbf{opinion raisonnée}, pas une
          certitude~; la présenter comme telle protège l'expert
          lors du contre-interrogatoire.
    \item La corrélation temporelle multi-sources exige une
          \textbf{synchronisation des horloges}~: documenter le
          décalage NTP de chaque système analysé avant toute chronologie.
\end{itemize}
\end{boxresume}

% ============================================================================
\section{L'Affaire Enron --- L'E-Discovery à l'Échelle (2001)}
\label{sec:enron}
% ============================================================================

\subsection{Contexte et artefact décisif}

La faillite d'Enron en décembre~2001 génère la première investigation
numérique à grande échelle de l'histoire~: 500~000~documents électroniques
(emails, feuilles Excel, présentations PowerPoint) doivent être collectés,
préservés, analysés et présentés en justice. L'enjeu dépasse le périmètre
technique~: pour la première fois, un tribunal doit statuer sur la
recevabilité de preuves numériques massives.

\medskip

\textbf{L'artefact décisif}~: l'\textbf{email d'Enron} du 23~août~2001,
dans lequel le vice-président Sherron Watkins avertit le PDG~: \textit{« I
am incredibly nervous that we will implode in a wave of accounting
scandals. »} Cet email n'était pas destiné à être conservé --- mais il
l'était, dans les sauvegardes automatiques des serveurs Exchange.

\subsection{Anatomie technique}

\subsubsection{La problématique de l'e-discovery à grande échelle}

500~000~documents ne peuvent pas être examinés manuellement. Enron a
conduit au développement des premières techniques de
\termecle{Technology-Assisted Review} (TAR)\index{TAR}, ancêtre des
outils d'IA forensique actuels. Les principes~:

\begin{enumerate}
    \item \textbf{Collecte préservante~:} extraction des emails depuis les
          serveurs Exchange avec préservation des métadonnées (horodatages,
          en-têtes SMTP, codes de priorité). La simple exportation en
          \texttt{.pst} aurait supprimé des informations cruciales.
    \item \textbf{Déduplication~:} les mêmes emails existent en copies
          multiples (boîtes envoi, reçu, sauvegardes). La déduplication
          basée sur le hachage permet de réduire le volume sans perte d'information.
    \item \textbf{Classification assistée~:} les premiers algorithmes de
          classification par mots-clés et par graphe de communication ont
          identifié les 200~personnes les plus impliquées parmi les
          20~000~boîtes email analysées.
\end{enumerate}

\subsubsection{La règle de l'email~: ce qui est envoyé ne disparaît jamais}

\begin{boxalert}
\textbf{Le principe de la persistance des communications électroniques~:}
un email envoyé existe en copies multiples --- dans la boîte d'envoi de
l'expéditeur, la boîte de réception du destinataire, les sauvegardes
automatiques des serveurs, les journaux des serveurs relais, les archives
légales si elles existent. La suppression d'une copie n'élimine pas les autres.

En contexte judiciaire~: la réquisition aux fournisseurs de services
(art.~30 Loi~2010/012 au Cameroun~; Convention de Budapest art.~17 et~18
pour la coopération internationale) est le mécanisme légal pour accéder
aux copies conservées par les opérateurs.
\end{boxalert}

\subsection{Leçon transférable}

\begin{boxresume}
\textbf{Enron --- Ce que vous emportez~:}
\begin{itemize}
    \item La collecte préservante des emails exige d'extraire les
          \textbf{métadonnées complètes} (en-têtes SMTP, horodatages
          de réception, Message-ID) --- pas seulement le corps du message.
    \item La déduplication basée sur le hachage est la méthode de référence
          pour traiter de grands volumes sans altérer les preuves.
    \item En droit camerounais~: la réquisition judiciaire aux opérateurs
          pour obtenir les communications conservées est l'article~30
          de la \loicam{2010/012}. Elle doit être formulée avec précision
          (période, type de données, critères d'identification).
\end{itemize}
\end{boxresume}

% ============================================================================
\section{Opération ONAMBELE (Cameroun)}
\label{sec:onambele}
% ============================================================================

\subsection{Contexte}

\begin{boxscenario}{Opération ONAMBELE --- Fraude MoMo}
\textbf{Situation initiale~:} Mme~ABENA Céline, infirmière à Yaoundé,
perd 850~000~FCFA en trois virements MTN MoMo (mars~2026) après avoir reçu
des SMS lui annonçant qu'elle a gagné une promotion MTN fictive. Les fonds
sont transférés vers le numéro de ZANGA Michel, 28~ans, qui a « prêté » son
compte à un individu connu sous le nom de « KOFFI » en échange de
5~000~FCFA. KOFFI est identifié comme KOUAME Narcisse, basé à Abidjan
(Côte d'Ivoire).
\end{boxscenario}

\subsection{Artefacts décisifs et anatomie technique}

L'affaire mobilise quatre artefacts complémentaires~:

\begin{boxartefact}{Historique de transactions MoMo}
\textbf{Source~:} réquisition judiciaire à MTN (\loicam{2010/012}, art.~30).

\textbf{Information probante~:} horodatages précis des trois virements,
numéro destinataire (ZANGA), et la cascade de transferts immédiats vers
trois autres numéros puis vers un retrait en espèces.

\textbf{Valeur forensique~:} la cascade (\textit{layering}) est un
schéma classique de blanchiment~; elle prouve que le numéro de ZANGA
n'était pas une destination finale mais un relais dans un réseau organisé.
\end{boxartefact}

\begin{boxartefact}{Métadonnées du check-in Facebook de KOFFI}
\textbf{Source~:} OSINT sur le profil Facebook de KOFFI.

\textbf{Information probante~:} un check-in dans un restaurant de Yaoundé
daté du 25~février~2026 --- deux semaines avant la fraude. Le profil indiquait
être à Abidjan, mais ce check-in contredit cette déclaration et place KOFFI
au Cameroun au moment des faits.

\textbf{Valeur forensique~:} les check-ins Facebook contiennent les
coordonnées GPS inscrites dans les métadonnées du post. En l'occurrence~:
latitude~3.86°N, longitude~11.52°E --- centre de Yaoundé.
\end{boxartefact}

\subsection{Erreurs à éviter}

\begin{boxalert}
\textbf{Erreur n°1 --- Arrêter à ZANGA~:} ZANGA est une « mule » --- un
intermédiaire recruté pour ouvrir un compte et le prêter. L'arrêter seul
ne démantèle pas le réseau. L'investigateur doit remonter la chaîne~:
ZANGA $\to$ KOFFI $\to$ le vrai organisateur.

\textbf{Erreur n°2 --- Négliger le téléphone allumé~:} le téléphone
Samsung de ZANGA était allumé lors de l'arrestation. Un OPJ non formé
pourrait le brancher sur un chargeur sans activer le mode avion. En
quelques secondes, un message de «~remote wipe~» pourrait effacer les
contacts de KOFFI sur WhatsApp. Procédure correcte~: mode avion
immédiat, puis ADB.

\textbf{Erreur n°3 --- Omettre la coopération internationale~:}
KOUAME Narcisse est en Côte d'Ivoire. Sans demande de coopération
(Convention de Budapest art.~29~; accord bilatéral Cameroun-CEDEAO),
il ne peut pas être arrêté. La demande doit partir dès l'identification,
pas attendre la fin de l'investigation nationale.
\end{boxalert}

\begin{boxCRO}
\textbf{Analyse CRO --- Opération ONAMBELE}

$\mathcal{C}$~: le contenu des messages WhatsApp entre ZANGA et KOFFI
est une preuve cruciale, mais son accès exige le déblocage du téléphone
(code PIN que ZANGA refuse de communiquer). Tenter un brute-force sans
mandat risque de violer $\mathcal{C}$ et compromettre $\mathcal{O}$.

$\mathcal{R}$~: les données MoMo (réquisition MTN) sont des preuves
fiables car issues directement des bases de l'opérateur, avec horodatages
serveur.

$\mathcal{O}$~: l'ensemble du dossier (réquisition MTN + OSINT Facebook
+ ADB après consentement ou mandat) est admissible si la procédure est
respectée.

\cro{$\sim$ à gérer}{$\uparrow$}{$\uparrow$ sous conditions}
\end{boxCRO}

% ============================================================================
\section{Opération TANGUE (Cameroun)}
\label{sec:tangue}
% ============================================================================

\subsection{Contexte}

\begin{boxscenario}{Opération TANGUE --- Forensique de disque}
BELLO Armand, comptable chez SONATRAV (entreprise de travaux publics),
est suspecté d'avoir détourné 25~millions FCFA via de faux bons de commande
créés dans le système comptable informatisé. Son ordinateur portable
Windows est saisi. Une image forensique E01 est créée avec FTK Imager.
Huit artefacts sont cachés dans l'image~: l'investigateur doit les trouver
tous avec Autopsy.
\end{boxscenario}

\subsection{Les huit artefacts et leur localisation}

\begin{tableaucours}{p{2.5cm}p{4.5cm}p{5cm}}{%
    \tableauHeader{Artefact} &
    \tableauHeader{Description} &
    \tableauHeader{Localisation Autopsy}
}
A1 --- Email virement &
    Email Outlook~: ordre de virement avec faux IBAN fournisseur &
    \texttt{Results → E-Mail Messages → rechercher 'virement'} \\
A2 --- Screenshot banque &
    Capture portail bancaire UBA après virement &
    \texttt{File Views → Images → filtrer par date} \\
A3 --- GPS dans photo &
    Photo prise à Yaoundé (pas à Douala comme déclaré) -- coordonnées GPS dans EXIF &
    \texttt{Results → EXIF Metadata → colonnes Latitude/Longitude} \\
A4 --- Document supprimé &
    \texttt{faux\_bon\_de\_commande.docx} marqué supprimé &
    \texttt{File Views → Deleted Files} \\
A5 --- Tableau détournements &
    Excel avec suivi des détournements~: montants, dates, bénéficiaires &
    \texttt{File Views → Recent Documents} \\
A6 --- Historique Chrome &
    Visite portail bancaire le jour du virement à 14h32 &
    \texttt{Results → Web History → filtrer par date} \\
A7 --- Prefetch keyfinder &
    \texttt{PASSWORDRECOVERY.EXE} -- outil d'extraction de mots de passe &
    \texttt{C:\textbackslash Windows\textbackslash Prefetch\textbackslash} \\
A8 --- Clé USB non déclarée &
    Kingston~16~Go, branchée le jour du virement à 14h40 &
    \texttt{Results → USB Device Attached} \\
\end{tableaucours}
\vspace{-0.8em}
\begin{center}{\small\itshape Tableau~3.1 --- Les huit artefacts de l'Opération
TANGUE et leur localisation dans Autopsy}\end{center}

\subsection{Leçon transférable}

L'affaire TANGUE illustre la \textbf{convergence probatoire}~: aucun des
huit artefacts n'est suffisant seul, mais leur combinaison est accablante.
L'email prouve l'intention~; le screenshot prouve l'action~; le GPS
invalide l'alibi~; le document supprimé prouve la préméditation~; le
Prefetch du keyfinder prouve la préparation technique~; la clé USB prouve
l'exfiltration.

\begin{boxCRO}
\textbf{Analyse CRO --- Opération TANGUE}

L'artefact A5 (tableau de suivi des détournements) pose une question
CRO intéressante~: un fichier \emph{auto-incriminant} a une valeur
probatoire réelle ($\mathcal{R}$ haute~: l'auteur lui-même documente
ses actes) mais doit être corroboré ($\mathcal{O}$ conditionnelle~: sans
corroboration, un juge prudent le considèrera comme insuffisant seul).

La corrélation A5 + A1 + A6 (même journée, même heure, même montant)
constitue la triade probatoire qui rend le dossier incontestable.
\cro{$\sim$}{$\uparrow\uparrow$}{$\uparrow$ par convergence}
\end{boxCRO}

% ============================================================================
\section{Opération MVOM (Cameroun)}
\label{sec:mvom}
% ============================================================================

\subsection{Contexte}

\begin{boxscenario}{Opération MVOM --- Ransomware sur PME}
TRANSBOIS SA, PME forestière de Yaoundé (45~employés, 2~milliards FCFA
de CA), subit une attaque ransomware LockBit. 85\% des fichiers sont chiffrés
(extension \texttt{.lockbitcm}). La rançon demandée est de 5~bitcoins
(\textasciitilde 400~000~USD). L'entrée initiale~: une connexion RDP
depuis une IP roumaine, utilisant le compte administrateur dont le mot de
passe était \texttt{Transbois2023}. 2,3~Go de données ont été exfiltrées
\emph{avant} le chiffrement.
\end{boxscenario}

\subsection{Anatomie technique~: reconstruction sur la Kill Chain}

La Kill Chain de Lockheed Martin en sept phases est l'outil de référence
pour reconstituer une attaque~:

\begin{tableaucours}{p{1.8cm}p{3.0cm}p{6.5cm}}{%
    \tableauHeader{Phase} &
    \tableauHeader{Nom} &
    \tableauHeader{Preuve dans l'affaire MVOM}
}
Phase~1 & Reconnaissance & Non documentée~; probable scan RDP depuis l'IP roumaine \\
Phase~2 & Armement & Préparation du ransomware LockBit~; hors périmètre camerounais \\
Phase~3 & Livraison & Connexion RDP avec \texttt{admin\_tb / Transbois2023} \\
Phase~4 & Exploitation & Accès administrateur obtenu (12/03~02h14) \\
Phase~5 & Installation & Dépôt de \texttt{enc\_all.ps1} dans \texttt{\%TEMP\%}~; processus \texttt{svchost.exe} suspect \\
Phase~6 & C\&C & Requête DNS vers \texttt{lockbit-pay-cm.onion}~; connexion vers port~4444 \\
Phase~7 & Actions & Exfiltration 2,3~Go (11/03) puis chiffrement (13/03) \\
\end{tableaucours}
\vspace{-0.8em}
\begin{center}{\small\itshape Tableau~3.2 --- Reconstruction MVOM sur la Kill Chain
(d'après l'analyse Wireshark + Volatility)}\end{center}

\subsection{L'analyse Volatility~: preuve de l'installation}

\begin{lstlisting}[language=bash_forensique,
    caption={Volatility 3 --- analyse mémoire RAM TRANSBOIS}]
# Lister les processus -- identifier le svchost.exe suspect
python3 vol.py -f transbois.mem windows.pslist
# Résultat : PID 4892 svchost.exe -- parent : cmd.exe (PID 4891)
# ANORMAL : svchost.exe ne devrait pas avoir cmd.exe comme parent

# Analyser les connexions réseau du processus suspect
python3 vol.py -f transbois.mem windows.netscan | grep 4892
# Résultat : connexion ESTABLISHED vers 185.220.XX.XX:4444

# Lister les fichiers ouverts par le processus
python3 vol.py -f transbois.mem windows.filescan | grep enc_all
# Résultat : \Device\HarddiskVolume3\Users\admin\AppData\Local\Temp\enc_all.ps1
\end{lstlisting}

\subsection{Leçon transférable~: la double victime}

\begin{boxalert}
\textbf{TRANSBOIS est simultanément victime et contrevenante.} Le fait
d'être victime d'une attaque ransomware n'exonère pas TRANSBOIS de ses
obligations légales~: absence de politique de mots de passe robuste (art.~27
\loicam{2024/017}), non-notification de violation de données à l'ANPDP
(art.~22 \loicam{2024/017}), absence de registre de traitement (art.~29
\loicam{2024/017}).

\textbf{Principe~:} un responsable de traitement qui facilite, même
involontairement, une attaque contre les données de ses clients engage
sa responsabilité administrative et potentiellement pénale. Les deux
procédures (victimisation + manquements) sont indépendantes et doivent
être traitées séparément.
\end{boxalert}

% ============================================================================
\section{Opération WOURI (Cameroun)}
\label{sec:wouri}
% ============================================================================

\subsection{Contexte}

\begin{boxscenario}{Opération WOURI --- Réseau BEC transfrontalier}
Un réseau de quatre individus opérant depuis trois pays (Cameroun,
Nigeria, Émirats arabes unis) défraude des entreprises camerounaises via des
attaques \termecle{BEC} (\textit{Business Email Compromise})~: usurpation
du domaine d'un fournisseur pour intercepter des virements. Deux victimes~:
SOTRADI SA (47M~FCFA) et AGROBOIS Cameroun (32M~FCFA). Rôles~:
ATANGANA (organisateur, Cameroun), OKAFOR (technicien, Nigeria),
HASSAN (blanchisseur, EAU), NGUEMO (complice-mule, Cameroun).
\end{boxscenario}

\subsection{Cartographie du réseau criminel}

\begin{figure}[H]
\centering
\begin{tikzpicture}[
    suspect/.style={
        draw=CouleurPrimaire, fill=CouleurDef,
        rounded corners=4pt, minimum width=3.2cm, minimum height=0.9cm,
        font=\small\bfseries\color{CouleurPrimaire}, align=center
    },
    victime/.style={
        draw=CouleurBordAlerte, fill=CouleurAlerte,
        rounded corners=4pt, minimum width=3.2cm, minimum height=0.9cm,
        font=\small\color{CouleurBordAlerte}, align=center
    },
    flux/.style={-{Stealth[length=5pt]}, thick, color=CouleurSecondaire},
    blanchiment/.style={-{Stealth[length=5pt]}, thick, dashed, color=CouleurAccent}
]
% Suspects
\node[suspect] (atangana) at (0, 0)    {ATANGANA\\Yaoundé (CM)};
\node[suspect] (okafor)   at (-4, -2)  {OKAFOR\\Lagos (NG)};
\node[suspect] (nguemo)   at (4,  -2)  {NGUEMO\\Yaoundé (CM)};
\node[suspect] (hassan)   at (0,  -4)  {HASSAN\\Dubaï (EAU)};
% Victimes
\node[victime] (sotradi)  at (-3.5, 2) {SOTRADI SA\\47M FCFA};
\node[victime] (agrobois) at (3.5,  2) {AGROBOIS\\32M FCFA};
% Flux
\draw[flux] (okafor)   -- node[left, font=\tiny, color=CouleurSecondaire]
    {phishing/BEC} (sotradi);
\draw[flux] (okafor)   -- node[right, font=\tiny, color=CouleurSecondaire]
    {phishing/BEC} (agrobois);
\draw[flux] (atangana) -- node[above, font=\tiny] {cibles} (okafor);
\draw[flux] (sotradi)  -- node[above left, font=\tiny, color=CouleurBordAlerte]
    {virement} (nguemo);
\draw[flux] (agrobois) -- node[above right, font=\tiny, color=CouleurBordAlerte]
    {virement} (nguemo);
\draw[blanchiment] (nguemo) -- node[right, font=\tiny, color=CouleurAccent]
    {BTC mixeur} (hassan);
\draw[blanchiment] (atangana) to[bend right=20]
    node[below, font=\tiny, color=CouleurAccent] {coordination} (hassan);
\end{tikzpicture}
\caption{Réseau criminel Opération WOURI --- flux d'information et financiers}
\end{figure}

\subsection{L'OSINT comme technique forensique à part entière}

L'identification d'OKAFOR (basé à Lagos, jamais entré au Cameroun) repose
exclusivement sur l'\termecle{OSINT}\index{OSINT} (\textit{Open Source
Intelligence})~:

\begin{enumerate}
    \item \textbf{LinkedIn~:} profil de « cybersecurity consultant »,
          3~anciens emplois vérifiables, email professionnel récupérable.
    \item \textbf{GitHub~:} un script de phishing publié avec des
          commentaires en anglais et en pidgin nigérian, signé du même
          pseudo.
    \item \textbf{Sherlock~:} l'outil de recherche de pseudo
          (\lstinline[language=bash_forensique]{sherlock duke_cyber_ng})
          identifie 12~plateformes supplémentaires.
    \item \textbf{Corrélation blockchain~:} les adresses Bitcoin
          transmises par ATANGANA sont reliées à des adresses connues du
          groupe SilverTerrier sur Blockchair.
\end{enumerate}

\begin{boxPNC}
\textbf{Procédure WOURI --- Actions coopération internationale}

L'affaire implique trois pays~; la coopération doit être lancée en parallèle
dès l'identification des suspects étrangers~:

\begin{enumerate}
    \item \textbf{OKAFOR (Nigeria)~:} transmission d'une Notice Bleue
          INTERPOL via le BCN Cameroun~; demande d'entraide judiciaire
          (ENTR) via l'ANIF si des flux financiers sont identifiés au Nigeria.

    \item \textbf{HASSAN (Émirats arabes unis)~:} demande de conservation
          urgente à Binance via le portail LEA (\textit{Law Enforcement
          Agency})~; demande d'identification du titulaire du compte Bitcoin.
          Les EAU ont signé la Convention de Budapest en 2021~: la
          procédure s'applique.

    \item \textbf{OKAFOR -- arrestation~:} une Notice Rouge INTERPOL est
          nécessaire pour permettre une arrestation provisoire au Nigeria.
          Elle requiert un mandat d'arrêt camerounais préalable.
\end{enumerate}
\end{boxPNC}

\begin{boxCRO}
\textbf{Analyse CRO --- Opération WOURI}

$\mathcal{C}$~: les communications WhatsApp d'ATANGANA (trouvées sur
son téléphone déverrouillé) sont une preuve directe mais contiennent
aussi des communications sans rapport avec l'affaire (famille, travail).
Le principe de proportionnalité impose de ne produire en justice que
les communications pertinentes.

$\mathcal{R}$~: les données blockchain (Blockchair) sont publiques et
reproductibles~; leur fiabilité est très haute. Les données Binance
(portail LEA) sont fiables mais leur chaîne d'obtention doit être
documentée pour $\mathcal{O}$.

$\mathcal{O}$~: le principal risque d'irrecevabilité est la coopération
internationale~: si les données obtenues du Nigeria ou des EAU ne
respectent pas les procédures \sigle{MLAT} (Mutual Legal Assistance
Treaty), un tribunal camerounais peut refuser de les utiliser.

\cro{$\downarrow$ gérer proportionnalité}{$\uparrow$}{$\sim$ dépend de la coopération}
\end{boxCRO}

% ============================================================================
\section{Tableau de Synthèse~: les Leçons Permanentes}
\label{sec:synthese-affaires}
% ============================================================================

Sept affaires, sept leçons, toutes applicables demain matin sur le terrain~:

\begin{tableaucours}{p{2.5cm}p{3.5cm}p{5.5cm}}{%
    \tableauHeader{Affaire} &
    \tableauHeader{Innovation forensique} &
    \tableauHeader{Leçon permanente}
}
BTK / Rader &
    Métadonnées comme preuve décisive &
    Extraire systématiquement les métadonnées~; une couche invisible peut
    identifier un suspect \\
Stuxnet &
    Cyberforensique industrielle &
    La sophistication est un marqueur d'attribution~;
    les ICS/SCADA sont une frontière émergente \\
WannaCry &
    Kill switch et kill chain &
    L'attribution est probabiliste~;
    la synchronisation des horloges est critique \\
Enron &
    E-discovery à grande échelle &
    La réquisition aux opérateurs est le levier légal du volume~;
    la déduplication par hachage est la méthode \\
ONAMBELE &
    Forensique Mobile Money + OSINT &
    Remonter la chaîne complète~;
    traiter le téléphone allumé selon le protocole avant tout \\
TANGUE &
    Convergence probatoire par artefacts multiples &
    Aucun artefact seul ne suffit~;
    la convergence de huit preuves est accablante \\
MVOM &
    Kill chain + Volatility + Loi 2024 &
    La victime peut être aussi contrevenante~;
    les deux procédures sont indépendantes \\
WOURI &
    BEC + OSINT + coopération multilatérale &
    Lancer la coopération internationale dès l'identification~;
    les plateformes \sigle{LEA} sont des outils légaux directement accessibles \\
\end{tableaucours}
\vspace{-0.8em}
\begin{center}{\small\itshape Tableau~3.3 --- Synthèse des leçons permanentes}\end{center}

\separateur

\begin{boxresume}
\textbf{Ce qu'il faut retenir de ce chapitre~:}

\begin{itemize}
    \item Chaque grande affaire a produit une \textbf{innovation forensique}
          qui est aujourd'hui une technique standard~: les métadonnées (BTK),
          le sandboxing comportemental (Stuxnet), l'analyse kill chain
          (WannaCry, MVOM), la réquisition numérique (Enron, ONAMBELE).

    \item L'\textbf{artefact décisif} --- la seule preuve qui a fait basculer
          le dossier --- est rarement la plus spectaculaire. C'est souvent
          une métadonnée oubliée, un horodatage incohérent, ou un processus
          système anormal.

    \item Les \textbf{erreurs récurrentes} --- brancher le téléphone sans
          mode avion, arrêter le maillon faible sans remonter la chaîne,
          négliger la coopération internationale --- sont évitables. Elles
          se préviennent par le protocole, pas par le talent.

    \item Les \textbf{cas camerounais} (ONAMBELE, TANGUE, MVOM, WOURI)
          illustrent que les schémas mondiaux (fraude MoMo, ransomware BEC,
          réseaux transfrontaliers) opèrent au Cameroun avec les mêmes
          mécanismes et la même sophistication qu'ailleurs.

    \item Le Trilemme CRO se manifeste différemment dans chaque affaire~:
          Stuxnet sacrifie $\mathcal{O}$ au profit de $\mathcal{C}$ et
          $\mathcal{R}$~; TANGUE maximise $\mathcal{R}$ et $\mathcal{O}$~;
          WOURI dépend de la coopération internationale pour atteindre
          $\mathcal{O}$.
\end{itemize}
\end{boxresume}

\bigskip

\begin{boxexo}[title={\ding{45}\quad Exercice~3.1 --- Identification de l'artefact décisif \niveaubase}]
Pour chacune des situations suivantes, identifiez l'artefact décisif probable
et la méthode qui permettrait de l'extraire~:

\begin{enumerate}[(a)]
    \item Un suspect nie avoir participé à une réunion en ligne, mais vous
          disposez de son ordinateur Windows.
    \item Une entreprise prétend que ses données n'ont pas été exfiltrées
          avant le chiffrement par un ransomware. Vous disposez d'une
          capture réseau Wireshark de 48~heures.
    \item Une escroquerie sentimentale via Facebook~: le suspect prétend
          être à Paris, mais la victime affirme avoir été contactée depuis
          le Cameroun. Vous disposez du profil Facebook du suspect.
\end{enumerate}
\end{boxexo}

\begin{boxexo}[title={\ding{45}\quad Exercice~3.2 --- Reconstitution d'attaque \niveaumoyen}]
\textbf{Scénario MVOM simplifié~:} vous recevez les données suivantes
concernant une PME camerounaise attaquée~:

\begin{itemize}
    \item Log pare-feu~: connexion RDP entrante depuis \texttt{91.192.X.X}
          le 03/04/2026 à 23h47, durée 3h12.
    \item Processus Volatility~: \texttt{powershell.exe} (PID~7231),
          parent~: \texttt{cmd.exe} (PID~7228), connexion vers
          \texttt{91.192.X.X:443} pendant 2h45.
    \item Log pare-feu~: transfert sortant de 4,7~Go vers \texttt{185.220.Y.Y}
          le 04/04/2026 entre 01h15 et 02h58.
    \item Autopsy~: 87\% des fichiers \texttt{.xlsx} et \texttt{.docx}
          portent l'extension \texttt{.locked2026}.

\end{itemize}

(a)~Placez chaque élément sur la Kill Chain (7~phases).\\
(b)~Quel est le vecteur d'entrée probable~? Quelle preuve l'établit~?\\
(c)~Quand a eu lieu l'exfiltration~: avant ou après le chiffrement~?
Quelle implication pour les victimes (données personnelles exfiltrées)~?\\
(d)~Quels articles de la \loicam{2010/012} et de la \loicam{2024/017}
qualifier pour les auteurs~? Pour la PME elle-même~?
\end{boxexo}

\begin{boxexo}[title={\ding{45}\quad Exercice~3.3 --- Analyse CRO d'une décision
investigative \niveauexpert}]
Dans l'Opération WOURI, l'OPJ en charge décide d'extraire l'intégralité
du téléphone d'ATANGANA via ADB (après que celui-ci a fourni son code PIN)
pour analyser toutes ses communications des 12~derniers mois.

\begin{enumerate}[(a)]
    \item Évaluez cette décision selon le Trilemme CRO~: quelles
          dimensions sont renforcées ou fragilisées~?
    \item La décision est-elle proportionnelle~? Argumentez en référence
          au cadre légal camerounais et au type d'infraction.
    \item Proposez une alternative moins intrusive qui préserverait mieux
          $\mathcal{C}$ tout en maintenant une $\mathcal{O}$ satisfaisante.
    \item Comment documentez-vous dans votre rapport le choix effectué
          et sa justification, pour qu'un juge puisse l'évaluer~?
\end{enumerate}
\end{boxexo}

\bigskip

\begin{boxinfo}
\textbf{Transition.} Les Chapitres~4 et~5 complètent la Partie~I en
formalisant les \textbf{fondements théoriques} mobilisés dans ces affaires
(Shannon, Dolev-Yao, théorie de la complexité) et en dressant l'\textbf{état
de l'art} 2025 (IA forensique, deepfakes, conteneurs). La Partie~II
construit ensuite le \textbf{processus investigatif universel} qui structure
la mise en œuvre de toutes ces leçons dans une investigation réelle.
\end{boxinfo}
